\scanpage{67}
% archive foliation: 66
De cette façon, pour $\underline W$ fixé, \deleted{\unsure{$\mathfrak X$ fixé}} \deleted{\unsure{fixé}} $\mathfrak X$ et $W'$ variables,
\[
W'^a(\mathfrak X\times_{W(S)}W'(S)),
\]
devient un \uline{bifoncteur} en $\mathfrak X$, $W'$, \uline{covariant} en $\mathfrak X$, \uline{contravariant en} $W'$.

\textcircled{2} \uline{Critère d'existence : Cas affine}
\margin{Ici, il suffit que $W$ \unsure{en plus} transforme affines en affines \deleted{\unsure{sur un}}}
\margin{Foncteurs $W$ \uline{spécial} [\unsure{associé à un} schéma \deleted{\unsure{d'}} en anneaux sur $S$, affine sur $S$]}

Supposons que $\underline W$ transforme schémas affines (sur ouverts affines de $S$) en schémas affines. Donc, au dessus d'un ouvert affine de $S$, d'anneau $A$, $\underline W$ est défini par un foncteur covariant, également noté $\underline W$, des \unsure{fais.} $A$-algèbres \inserted{$B$} en anneaux $\underline W(B)$. [\unsure{Avec} l'hypothèse de \deleted{localité} \inserted{\unsure{caractère local}} de ce foncteur, que nous n'expliciterons pas]. On suppose qu'on a
\inserted{$\underline W(B)=\mathrm{Hom}_{A\text{-alg.}}(\Lambda,B)$}
où $\Lambda$ est \deleted{l'algèbre affine} \inserted{l'anneau} d'un schéma en anneaux \deleted{affine} sur $S$, affine sur $S$.
